%% kapa: exam=1 round=20 year=2009 period=1 form=A subject=economics subject_ko=경제원론 seq=30 number=70
%% source: Q-Net 감정평가사 기출문제 (cst00309) artlSeq=5024178 file=extracted/1008/5024178/2009년 제20회 감정평가사 1차시험 기출문제/1교시.hwp
%% conversion: kordoc 4.24.51 -> tools/gen1.py (auto); review=proofread-sonnet; flags=tikz
\begin{question}{70}
다음 그림에 관한 설명으로 옳지 않은 것은? (단, $I_{ 1 }$ 과 $I_{ 2 }$ 는 일반적인 특성을 갖는 무차별곡선이고, 우하향하는 선분은 예산선을 나타낸다.)
\begin{center}
\begin{tikzpicture}[font=\small, x=1cm, y=1cm]
  \draw (0,0) -- (0,5.2);
  \draw (0,0) -- (5.6,0);
  \node[above] at (0,5.2) {Y재};
  \node[below] at (5.3,-0.1) {X재};
  \node[below left] at (0,0) {0};
  \draw (0,3.98) -- (4.35,0);
  \draw plot[smooth] coordinates {(0.69,4.38) (0.76,3.8) (0.84,3.23) (1.27,1.93) (1.97,1.23) (3.49,0.78) (4.94,0.63)};
  \draw plot[smooth] coordinates {(1.6,4.78) (1.82,4.05) (2.27,2.95) (2.97,2.25) (3.97,1.85) (5.34,1.28)};
  \fill (0.84,3.23) circle (2.2pt);
  \fill (3.49,0.78) circle (2.2pt);
  \node at (0.5,3.05) {A};
  \node at (3.4,0.45) {B};
  \node at (0.85,4.1) {$I_1$};
  \node at (1.9,4.5) {$I_2$};
\end{tikzpicture}
\end{center}
\begin{choices}
  \item Y재로 표시한 X재의 한계대체율이 B점보다 A점에서 크다.
  \item 무차별곡선 $I_{ 1 }$ 에서의 상품묶음이 $I_{ 2 }$ 에서의 어떤 상품묶음보다도 효용이 작다.
  \item 소비자가 A점에서 얻는 총효용의 크기가 B점에서 얻는 총효용의 크기와 같다.
  \item A점에서 X재의 1원당 한계효용은 Y재의 1원당 한계효용보다 작다.
  \item B점에서 소비하는 경우, 효용을 극대화하기 위해서는 X재의 소비를 감소시키고 Y재의 소비를 증가시켜야 한다.
\end{choices}
\end{question}
\answer{4}
